% =====================================================================
\section{Validación de requisitos}
\label{sec:validacion}
 
\subsection{Walkthrough estructurado en lugar de inspección formal}
 
\begin{table}[H]
\centering\footnotesize
\caption{Parámetros de la validación aplicada al ERS de SGCV-IA}
\label{tab:validacion}
\begin{tabular}{|L{6.4cm}|L{8.6cm}|}
\hline
\rowcolor{grisclaro}\textbf{Parámetro} & \textbf{Valor}\\ \hline
Documento validado & ERS/SRS de SGCV-IA, versiones 2.0 y 3.0 (Entrega 2A)\\ \hline
Técnica & \emph{Walkthrough} estructurado, integrado al cierre de cada iteración de Scrum\\ \hline
Alcance & Documento completo, sin exclusión declarada de apéndices\\ \hline
Registro de hallazgos & No se llevó un registro formal de defectos con severidad y tipo; las correcciones se aplicaron directamente sobre las fuentes en cada iteración\\ \hline
Evidencia & Carpeta \texttt{02\_Evidencias/Validacion\_Walkthrough/} del repositorio\\ \hline
\end{tabular}
\end{table}
 
\noindent
SGCV-IA no aplicó el método de Fagan con roles diferenciados, preparación individual obligatoria y
prohibición de discutir soluciones en sesión \cite{fagan1976}. La decisión y su justificación se
desarrollan en la \S\ref{sec:metodologia}: un ciclo de inspección formal completo exige una cadencia
de reunión y una preparación previa incompatibles con \emph{sprints} cortos y un equipo reducido. La
consecuencia directa de esa decisión es metodológica y conviene declararla sin rodeos: al no existir
una lista de verificación de inspección ni una consolidación de defectos por severidad y tipo, este
documento no puede presentar las métricas de densidad de defectos, distribución por tipo o esfuerzo
en horas-persona que sí serían exigibles tras una inspección formal tipo Fagan. Lo que sí puede
documentarse, con evidencia trazable al historial de versiones del ERS, son las correcciones
concretas que el \emph{walkthrough} produjo entre la versión 2.0 y la versión 3.0 del documento.
 
\subsection{Correcciones producidas por el walkthrough}
 
\begin{table}[H]
\centering\footnotesize
\caption{Correcciones aplicadas al ERS entre las versiones 2.0 y 3.0}
\label{tab:correcciones}
\begin{tabular}{|C{2.2cm}|L{6.0cm}|L{6.4cm}|}
\hline
\rowcolor{grisclaro}\textbf{Elemento} & \textbf{Estado antes de la corrección} & \textbf{Corrección aplicada}\\ \hline
Numeración de RF & \id{RC-22} (``Bajo costo de implementación y operación'') ocupaba una fila de requisito funcional & Reclasificado como restricción de diseño (\id{RST-01}); dejó de ocupar fila de RF\\ \hline
Identificador \id{RF-25} & Aprobado previamente por el CCB en una práctica anterior para un requisito de contenido distinto (asistencia a citas) & Reservado y no reutilizado; el requisito de portal de derechos se formalizó como \id{RF-26} para evitar colisión de identificadores\\ \hline
Requisitos no funcionales & Ausentes de la Entrega 2A inicial & Incorporados en la v3.0: 18 RNF (\id{RNF-01} a \id{RNF-18}), superando el mínimo de 15 exigido por la rúbrica\\ \hline
Consistencia y formato & Referencias internas y consistencia de formato pendientes de revisión final & Corregidas y declaradas explícitamente en el historial de cambios de la v3.0\\ \hline
\end{tabular}
\end{table}
 
\noindent
El patrón de estas correcciones es distinto del que produce una inspección formal por severidad y
tipo, pero es igualmente informativo: las dos primeras filas son de consistencia de identificadores
---un requisito ocupando la categoría equivocada, un identificador reutilizado antes de tiempo--- y
no de contenido faltante; la tercera es de completitud del catálogo, no de un elemento individual mal
especificado. Es el mismo tipo de defecto que la \S\ref{sec:metodologia} anticipa como punto débil de
una lista de verificación centrada en revisar elementos uno por uno: nada obliga al revisor a
preguntarse si el conjunto está completo o si la numeración es internamente coherente, y ambos
problemas aparecieron aquí.
 
\subsection{Por qué ocurrieron: análisis de causa}
 
Las cuatro correcciones de la tabla anterior no son independientes entre sí; comparten dos causas.
 
\textbf{Las dos primeras son de gobierno de identificadores entre entregas.} \id{RC-22} ocupó una
fila de requisito funcional porque en el momento de elicitarlo no existía todavía la categoría de
restricción de diseño en el catálogo; la reclasificación llegó cuando esa categoría se formalizó, no
antes. El caso de \id{RF-25}/\id{RF-26} tiene una causa distinta y más estructural: el proyecto
reutiliza un espacio de identificadores entre prácticas experimentales distintas del mismo curso
(PE4\_U4 y esta entrega), y ese espacio no está centralizado en un único documento de control. El
identificador RF-25 ya estaba baselineado por el Change Control Board para un requisito de contenido
distinto (``Confirmación de asistencia a citas'') antes de que este ERS llegara a numerar el portal
de derechos; el conflicto se detectó porque el equipo revisó el histórico antes de baselinear, no
porque un procedimiento lo hubiera bloqueado automáticamente.
 
\textbf{Las dos últimas son de completitud del catálogo, no de corrección de un elemento.} La
incorporación tardía de los 18 RNF y las correcciones de consistencia declaradas en la v3.0 responden
a la misma limitación que la \S\ref{sec:validacion} ya señaló para el \emph{walkthrough}: revisar si
cada elemento presente está bien escrito no obliga a preguntarse si faltan elementos completos por
declarar. El catálogo de RNF no estaba mal especificado en la Entrega 2A inicial; sencillamente no
estaba.
 
\subsection{El precedente del Change Control Board}
 
SGCV-IA sí cuenta con un precedente real de gobierno formal de cambios: el caso \id{RF-25}/\id{RF-26}
muestra que existe un Change Control Board que baselinea identificadores entre prácticas del curso, y
que ese baseline se respeta incluso cuando resulta incómodo ---obligó a renumerar un requisito ya
redactado en este ERS en lugar de reabrir el identificador ya aprobado. Lo que no existe para esta
entrega es un CCB actuando \emph{sobre} este ERS a partir de hallazgos de una inspección formal: las
tres correcciones restantes de la tabla anterior se resolvieron por consenso del equipo durante el
\emph{walkthrough}, sin tramitarse como solicitud de cambio formal, precisamente porque no derivaban
de una inspección que las hubiera registrado como defectos con ID propio.
 
\subsection{Limitación declarada}
 
Conviene ser explícito sobre lo que este documento no puede afirmar. Sin un registro de defectos
consolidado por severidad, tipo y esfuerzo, no es posible calcular una densidad de defectos por
página ni una tasa de detección por rol, y este ERS no las reporta ni las estima. Afirmar cifras de
ese tipo sin la inspección que las produce sería presentar como medición lo que sería una suposición.
 
\subsection{Acción correctiva de proceso}
 
La conclusión operativa no es que el \emph{walkthrough} haya sido insuficiente, sino que puede
complementarse sin adoptar el costo completo de una inspección formal. Para la siguiente entrega se
recomiendan dos ajustes concretos: mantener una lista de verificación mínima que registre cada
corrección aplicada durante el \emph{walkthrough} ---qué cambió, en qué sección, por qué--- en lugar
de aplicar la corrección directamente sin dejar rastro; e incluir explícitamente en cada revisión los
apéndices y los conteos globales del documento cuando este haya sufrido cambios aprobados, que es
precisamente donde se originaron los dos primeros hallazgos de la tabla anterior. Ese registro
mínimo no exige los roles ni la cadencia de Fagan, pero sí permitiría, en una futura entrega, calcular
métricas de proceso que hoy este documento no tiene base para reportar.
 
